\section{Introduction}
\label{sec:intro}

An AI system that behaves as intended during evaluation may behave differently when the conditions of supervision change. Alignment faking is a particularly consequential version of this problem: a model may comply with a training objective because it expects noncompliance to modify its future behavior, rather than because it endorses that objective \citep{greenblatt_alignment_faking}. Such behavior complicates the interpretation of safety evaluations, since apparent compliance need not demonstrate stable alignment.

Evidence of alignment faking in a large model does not establish how broadly the phenomenon occurs at smaller scales. \citet{koorndijk_empirical_2025} reported a prompt-induced compliance gap in LLaMA 3.1 8B and found that scratchpad prompting reduced the observed signal. Replication across model families and behavioral scenarios is needed to determine whether this result extends beyond the original model and context.

We investigate this question with locally served, quantized language models. The experimental intervention is textual: models are told that free-tier conversations will be used for retraining and paid-tier conversations will not. No model is retrained during the experiments. We compare responses under these descriptions across three scenarios and two reasoning conditions. Our scope is therefore behavior under simulated training incentives, rather than deception demonstrated during an actual training process.

The main contributions are:
\begin{enumerate}
  \item A replication of the harmful-request evaluation and an extension to 14 model configurations from five families under two scratchpad conditions.
  \item Two additional behavioral evaluations, persona preservation and animal welfare dismissal, that probe whether the compliance-gap pattern transfers across preference conflicts.
  \item A comparison of classifier-judged behavior, self-reported acceptance, and scratchpad classifications, with a limited human validation sample and manual review of classifier-positive scratchpads.
\end{enumerate}

The motivating literature review is summarized briefly in Section~\ref{sec:slr}. Most of the article concerns the experimental methods, results, and limits of interpretation. The central finding is a reproducible, context-specific compliance gap, rather than evidence that alignment faking is widespread in small language models.

